\chapter{Distribution}
\section{Project}
    the most standard structure compatible with PEP517 is:
    \dirtree{%
    .1 packagename\_project.
    .2 src.
    .3 packagename.
    .4 \_\_main\_\_.py.
    .4 \_\_init\_\_.py.
    .2 tests.
    .3 integration
    .3 performance
    .3 unit
    .3 performance
    .3 integration
    .2 dist.
    .2 .gitignore.
    .2 pyproject.toml.
    .2 README.md.
    }
    in unit, write unit tests, in performance write the benchz marks and intergration write data pipe line test
    
\section{Package}
    \paragraph{install/uninstall}
        navigate to the target folder in terminal while you are in the virtual environment and run
        \begin{minted}{shell}
            pip install -e .
        \end{minted}
        and to uninstall simple run the following anywhere in terminal
        \begin{minted}{shell}
            pip unintall $package_name
        \end{minted}

\section{Pypi}

\section{wheel}
    A wheel is a built package format — basically a pre-built, ready-to-install archive for a Python project.
    File extension: .whl
    It’s faster to install than building from source because all the compiling is already done.
    
    Wheels follow the PEP 427 standard.
    
    How to build:
    
    \# do all the follwing commands inside your venv
    
    \begin{lstlisting}[language=bash]
        cd ~/repos/utilityx_project
        pip install -U build
        #creates dist/.tar.gz and dist/.whl all go into the project build forlder
        python -m build
    \end{lstlisting}
    
    How to install
    
    \begin{lstlisting}[language=bash]
    pip install dist/utilityx-1.0.0-*.whl
    \end{lstlisting}